% GENERATED FILE. Edit canonical lessons, then run npm run generate:books.
% canonical-source-hash: fnv1a64:a691ed2359100d01
% canonical-lessons: TE-C546-curu, TE-C546-pita, TE-C546-bavi, TE-C546-cembu, TE-C546-boddinka

\chapter{Around the House}
\label{ch:te-around-the-house}
\hlchaptermodality{546}

\begin{chapteropening}
\textbf{By the end of this chapter:} \emph{I can name the eaves, a low stool, a well, a water pot and a cockroach.}

Complete the last lesson of chapter 546: I can name the eaves, a low stool, a well, a water pot and a cockroach.
\end{chapteropening}

\section[\texorpdfstring{\te{చూరు} (cūru)}{చూరు (cūru)}]{\te{చూరు} (cūru) — the eaves}

\label{lesson:TE-C546-curu}



\begin{quote}
Before the new one: say the Telugu for to shove, then the Telugu for a bet.
\end{quote}

\subsection*{You'll want to know: \te{చూరు}}
\textbf{\te{చూరు}} — \emph{cūru} — \textquotedblleft{}the eaves\textquotedblright{}.

\begin{quote}
\te{వర్షం} \te{నీళ్ళు} \te{చూరు} \te{నుంచి} \te{కారుతున్నాయి}.
\end{quote}

\begin{quote}
\emph{varṣaṁ nīḷḷu cūru nuñci kārutunnāyi.}
\end{quote}

\begin{quote}
Rainwater is dripping from the eaves.
\end{quote}

The edge of a sloping roof.

\begin{cousinweb}
\textbf{\te{చూరు}} (\emph{cūru}) is native Dravidian: DEDR 2729 lists Telugu \emph{cūru}, \textquotedblleft{}eaves, the edge of a roof\textquotedblright{}, with Kannada \emph{cūru}.
\end{cousinweb}

\begin{grammarlens}[title={What you've built}]
One more word for describing what you lived and felt.
\end{grammarlens}

\subsection*{Your turn}
\begin{itemize}
\raggedright
  \item \emph{Say it:} \emph{cūru}
  \item \emph{Say it:} \emph{cūru}, once more
  \item \emph{Say it:} \emph{varṣaṁ nīḷḷu cūru nuñci kārutunnāyi.} — Rainwater is dripping from the eaves.
  \item \emph{Recall:} say the Telugu for to shove, then the Telugu for a bet, then say \emph{cūru} again
\end{itemize}

\begin{tcolorbox}[breakable,colback=teal!4,colframe=teal!35!black,title={Before you move on}]
What does \te{చూరు} mean? (\textquotedblleft{}The eaves\textquotedblright{}.) Say it once more.
\end{tcolorbox}
\section[\texorpdfstring{\te{పీట} (pīṭa)}{పీట (pīṭa)}]{\te{పీట} (pīṭa) — a low wooden stool}

\label{lesson:TE-C546-pita}



\begin{quote}
Before the new one: say the Telugu for a bet, then the Telugu for the eaves.
\end{quote}

\subsection*{You'll want to know: \te{పీట}}
\textbf{\te{పీట}} — \emph{pīṭa} — \textquotedblleft{}a low wooden stool\textquotedblright{}.

\begin{quote}
\te{అమ్మమ్మ} \te{పీట} \te{మీద} \te{కూర్చుని} \te{కూరగాయలు} \te{కోసింది}.
\end{quote}

\begin{quote}
\emph{ammamma pīṭa mīda kūrcuni kūragāyalu kōsindi.}
\end{quote}

\begin{quote}
Grandma sat on a low stool and cut the vegetables.
\end{quote}

The low plank seat for eating or working on the floor.

\begin{cousinweb}
\textbf{\te{పీట}} (\emph{pīṭa}) is borrowed from Sanskrit \emph{pīṭha}, \textquotedblleft{}a stool, a bench\textquotedblright{} (CDIAL 8222).
\end{cousinweb}

\begin{grammarlens}[title={What you've built}]
One more word for describing what you lived and felt.
\end{grammarlens}

\subsection*{Your turn}
\begin{itemize}
\raggedright
  \item \emph{Say it:} \emph{pīṭa}
  \item \emph{Say it:} \emph{pīṭa}, once more
  \item \emph{Say it:} \emph{ammamma pīṭa mīda kūrcuni kūragāyalu kōsindi.} — Grandma sat on a low stool and cut the vegetables.
  \item \emph{Recall:} say the Telugu for a bet, then the Telugu for the eaves, then say \emph{pīṭa} again
\end{itemize}

\begin{tcolorbox}[breakable,colback=teal!4,colframe=teal!35!black,title={Before you move on}]
What does \te{పీట} mean? (\textquotedblleft{}A low wooden stool\textquotedblright{}.) Say it once more.
\end{tcolorbox}
\section[\texorpdfstring{\te{బావి} (bāvi)}{బావి (bāvi)}]{\te{బావి} (bāvi) — a dug well}

\label{lesson:TE-C546-bavi}



\begin{quote}
Before the new one: say the Telugu for the eaves, then the Telugu for a low wooden stool.
\end{quote}

\subsection*{You'll want to know: \te{బావి}}
\textbf{\te{బావి}} — \emph{bāvi} — \textquotedblleft{}a dug well\textquotedblright{}.

\begin{quote}
\te{మా} \te{పెరట్లో} \te{ఒక} \te{పాత} \te{బావి} \te{ఉంది}.
\end{quote}

\begin{quote}
\emph{mā peraṭlō oka pāta bāvi undi.}
\end{quote}

\begin{quote}
There is an old well in our backyard.
\end{quote}

\te{బావి} \te{నీళ్ళు}, well water.

\begin{cousinweb}
\textbf{\te{బావి}} (\emph{bāvi}) is usually traced to Sanskrit \emph{vāpī}, \textquotedblleft{}a pond, a tank\textquotedblright{} (CDIAL 11529).
\end{cousinweb}

\begin{grammarlens}[title={What you've built}]
One more word for describing what you lived and felt.
\end{grammarlens}

\subsection*{Your turn}
\begin{itemize}
\raggedright
  \item \emph{Say it:} \emph{bāvi}
  \item \emph{Say it:} \emph{bāvi}, once more
  \item \emph{Say it:} \emph{mā peraṭlō oka pāta bāvi undi.} — There is an old well in our backyard.
  \item \emph{Recall:} say the Telugu for the eaves, then the Telugu for a low wooden stool, then say \emph{bāvi} again
\end{itemize}

\begin{tcolorbox}[breakable,colback=teal!4,colframe=teal!35!black,title={Before you move on}]
What does \te{బావి} mean? (\textquotedblleft{}A dug well\textquotedblright{}.) Say it once more.
\end{tcolorbox}
\section[\texorpdfstring{\te{చెంబు} (cembu)}{చెంబు (cembu)}]{\te{చెంబు} (cembu) — a small water pot}

\label{lesson:TE-C546-cembu}



\begin{quote}
Before the new one: say the Telugu for a low wooden stool, then the Telugu for a dug well.
\end{quote}

\subsection*{You'll want to know: \te{చెంబు}}
\textbf{\te{చెంబు}} — \emph{cembu} — \textquotedblleft{}a small water pot\textquotedblright{}.

\begin{quote}
\te{తాతయ్య} \te{చెంబుతో} \te{నీళ్ళు} \te{తాగారు}.
\end{quote}

\begin{quote}
\emph{tātayya cembutō nīḷḷu tāgāru.}
\end{quote}

\begin{quote}
Grandpa drank water from a small pot.
\end{quote}

A small round metal pot with a neck.

\begin{cousinweb}
\textbf{\te{చెంబు}} (\emph{cembu}) is native Dravidian: DEDR 2775 lists Telugu \emph{cembu}, \textquotedblleft{}goblet, pitcher\textquotedblright{}, with Tamil \emph{cempu}, \textquotedblleft{}copper; a metal vessel\textquotedblright{}; DEDR also raises the possibility that Sanskrit \emph{tāmra}, \textquotedblleft{}copper\textquotedblright{}, played a part.
\end{cousinweb}

\begin{grammarlens}[title={What you've built}]
One more word for describing what you lived and felt.
\end{grammarlens}

\subsection*{Your turn}
\begin{itemize}
\raggedright
  \item \emph{Say it:} \emph{cembu}
  \item \emph{Say it:} \emph{cembu}, once more
  \item \emph{Say it:} \emph{tātayya cembutō nīḷḷu tāgāru.} — Grandpa drank water from a small pot.
  \item \emph{Recall:} say the Telugu for a low wooden stool, then the Telugu for a dug well, then say \emph{cembu} again
\end{itemize}

\begin{tcolorbox}[breakable,colback=teal!4,colframe=teal!35!black,title={Before you move on}]
What does \te{చెంబు} mean? (\textquotedblleft{}A small water pot\textquotedblright{}.) Say it once more.
\end{tcolorbox}
\section[\texorpdfstring{\te{బొద్దింక} (boddiṅka)}{బొద్దింక (boddiṅka)}]{\te{బొద్దింక} (boddiṅka) — a cockroach}

\label{lesson:TE-C546-boddinka}



\begin{quote}
Before the new one: say the Telugu for a dug well, then the Telugu for a small water pot.
\end{quote}

\subsection*{You'll want to know: \te{బొద్దింక}}
\textbf{\te{బొద్దింక}} — \emph{boddiṅka} — \textquotedblleft{}a cockroach\textquotedblright{}.

\begin{quote}
\te{వంటగదిలో} \te{బొద్దింకలు} \te{తిరుగుతున్నాయి}.
\end{quote}

\begin{quote}
\emph{vaṇṭagadilō boddiṅkalu tirugutunnāyi.}
\end{quote}

\begin{quote}
Cockroaches are running about in the kitchen.
\end{quote}

Plural: \te{బొద్దింకలు}, \emph{boddiṅkalu}.

Plural: \te{బొద్దింకలు}.

\begin{cousinweb}
\textbf{\te{బొద్దింక}} (\emph{boddiṅka}) is a Telugu word for a cockroach; its origin is not settled.
\end{cousinweb}

\begin{grammarlens}[title={What you've built}]
One more word for describing what you lived and felt.
\end{grammarlens}

\subsection*{Your turn}
\begin{itemize}
\raggedright
  \item \emph{Say it:} \emph{boddiṅka}
  \item \emph{Say it:} \emph{boddiṅka}, once more
  \item \emph{Say it:} \emph{vaṇṭagadilō boddiṅkalu tirugutunnāyi.} — Cockroaches are running about in the kitchen.
  \item \emph{Recall:} say the Telugu for a dug well, then the Telugu for a small water pot, then say \emph{boddiṅka} again
\end{itemize}

\begin{tcolorbox}[breakable,colback=teal!4,colframe=teal!35!black,title={Before you move on}]
What does \te{బొద్దింక} mean? (\textquotedblleft{}A cockroach\textquotedblright{}.) Say it once more.
\end{tcolorbox}
